%% 01-pendahuluan.tex — bab dokumen contoh proyek-laporan (hal. 29-31 (proposal), hal. 31-33 (laporan)).
\chapter{PENDAHULUAN}
\section{Gambaran Umum Masyarakat Sasaran}
Masyarakat sasaran beserta kondisi sosial ekonominya diuraikan pada subbab ini.

Kajian pustaka disusun dari sumber terkini \citep{creswell2018research, patten2017understanding}
dan dari pedoman resmi universitas \citep{unnes2024panduanakademik}. Regulasi penjaminan mutu
menjadi dasar kebijakan \citep{permendikbudristek2023}.
 Contoh rujukan gaya MLA: \citemla[23]{patten2017understanding}.

\section{Permasalahan}
Permasalahan yang dihadapi masyarakat sasaran diuraikan berdasarkan hasil observasi awal.

\section{Solusi Pemecahan Masalah}
Solusi yang ditawarkan disertai keunggulan dan kelemahannya. Mitra kegiatan sepakat
berkolaborasi seperti tertuang pada lampiran.

\section{Tujuan Proyek}
\section{Kontribusi Proyek bagi Masyarakat Sasaran}

\begin{table}[htbp]
  \centering
  \caption{Ringkasan penelitian atau kegiatan terdahulu yang relevan}
  \label{tab:terdahulu}
  \begin{tabular}{clc}
    \toprule
    No & Fokus kajian & Hasil utama \\
    \midrule
    1 & Kajian pertama \citep{martin2014write} & Efektif pada ranah kognitif \\
    2 & Kajian kedua \citep{geraldi2017project} & Meningkat kategori sedang \\
    \bottomrule
  \end{tabular}
  \sumber{penulis dari berbagai sumber (2026)}
\end{table}
